\documentclass[10pt,a4paper]{amsart}
\usepackage[margin=19mm]{geometry}
\usepackage{amsmath,amssymb,mathtools}
\usepackage[scr=rsfs,cal=euler]{mathalfa}
\usepackage{xfrac}
\usepackage[unicode,hidelinks,bookmarks=false]{hyperref}
\newcommand{\ie}{i.e.\ }
\newcommand{\cf}{cf.\ }
\newcommand{\resp}{respectively\ }
\newcommand{\Subsection}{\subsection}
\providecommand{\nobreakdash}{\nobreak\hbox{-}}
\setlength{\emergencystretch}{2em}
\multlinegap0pt
\usepackage{bm}
\usepackage{bbm}
\usepackage{dsfont}
\usepackage{xspace}
\usepackage{cancel}
\usepackage{mathtools}
\usepackage{amsthm}
\makeatletter
\@ifundefined{theo}{\newtheorem{theo}{Theo}}{}
\@ifundefined{lemme}{\newtheorem{lemme}{Lemma}}{}
\@ifundefined{prop}{\newtheorem{prop}[theo]{Proposition}}{}
\makeatother
\def\ep{{\varepsilon}}         
\def\ker{{\mathop{\rm Ker}}}
\newcommand{\N}{\mathbb{N}}
\newcommand{\R}{\mathbb{R}}
\newcommand{\ve}{\varepsilon}
\newcommand{\vp}{{\rm vp}}
\title[Extremising eigenvalues of the GJMS operators in a fixed con]{Extremising eigenvalues of the GJMS operators in a fixed conformal class}
\author{Emmanuel Humbert, Romain Petrides, Bruno Premoselli}
\date{Source: arXiv:2505.08280v2 (CC BY 4.0)}
\begin{document}
\maketitle

\noindent\emph{Source and licence.} Extremising eigenvalues of the GJMS operators in a fixed conformal class. Version arXiv:2505.08280v2 (CC BY 4.0), distributed under the Creative Commons Attribution 4.0 International licence, as recorded in the arXiv licence metadata for this exact version and shown on the abstract page https://arxiv.org/abs/2505.08280v2. Full rights notice: licenses/arxiv-2505.08280-CC-BY-4.0.txt.\par\smallskip

\noindent\textit{Editorial scope.} Counted unit: the Prop whose source label is \texttt{prop:wellposedmu\_1} in the article (source0.tex lines 894--911, source label \texttt{prop:wellposedmu\_1}), delivered together with the article's own complete original proof of it (source0.tex lines 1135--1264, one \texttt{proof} environment of 130 lines). No mathematics is written, repaired or re-derived: statement and proof are the article's own text line for line. The proof is detached from the statement in the source: the article states the result at lines 894--911 and prints its proof at lines 1135--1264, and the proof environment's own header names the statement, so the association is the article's own. Required context retained uncounted: eq:unique:continuation (lines 258--263); def:Kbeta (lines 494--496); eq:def:Ag (lines 595--597); lem:mainlemma (lines 705--719); eq:def:lambdak (lines 876--878); prop:uppersemicontinuity (lines 1018--1021); eq:enplus:kplus (lines 1103--1106); eq:enplus:kmoins (lines 1107--1109); prop:deficontinuityeigen (lines 1112--1132). Every definition, display and label of those lines is retained unchanged. Omitted from this package and not redistributed: 1--257, 264--493, 497--594, 598--704, 720--875, 879--893, 912--1017, 1022--1102, 1110--1111, 1133--1134, 1265--3559. Nothing inside a retained excerpt is omitted or altered: every retained body is delivered exactly as the article's lines are written, so no omission receipt of any kind accompanies this package. Citations: the retained excerpts carry no citation command at all, so no bibliography entry of the article is transcribed and no reference is redistributed. Reference adaptation: none was needed and none was made; every reference inside the delivered unit resolves inside it, so no unresolved reference or citation is produced. Printed slips kept literal: no printed word, symbol, hypothesis or number of the counted statement or of its proof was altered. Layout and class: the article's own class is \texttt{amsart} with packages including amsmath, amsrefs, amssymb, xcolor, hyperref, graphicx, dsfont; the wrapper is the lane's pinned \texttt{amsart} editorial wrapper at 10pt on A4, which restates the theorem-like environments the retained text needs and adds the article's own macro definitions that the retained text needs (\texttt{\textbackslash N}, \texttt{\textbackslash R}, \texttt{\textbackslash ep}, \texttt{\textbackslash ker}, \texttt{\textbackslash ve}, \texttt{\textbackslash vp}), with extra wrapper packages (bm, bbm, dsfont, xspace, cancel, mathtools). The delivered numbering is the unit's own. Assets: no retained span contains a figure environment, a graphics-inclusion command, a PDF-page inclusion, a table or a TikZ picture; no image file is copied into this package, the package declares no asset dependency and no image file is redistributed with it. Licence: version arXiv:2505.08280v2 of this article is distributed under the Creative Commons Attribution 4.0 International licence, as recorded in the arXiv licence metadata for this exact version and shown on the abstract page https://arxiv.org/abs/2505.08280v2. A full rights notice is delivered at licenses/arxiv-2505.08280-CC-BY-4.0.txt. Characters: every retained excerpt is pure ASCII, so the delivered engine, XeLaTeX with its default Latin Modern text font, reports no missing character.
\begin{equation} \label{eq:unique:continuation} 
\begin{aligned}
\text{ If } & \vp \in C^\infty(M)  \text{ satisfies }  P_g^s \vp = 0 \text{ and vanishes on a set }  \\
& \qquad \qquad \text{ of positive measure, then } \vp = 0. 
\end{aligned}
\end{equation}

\begin{equation} \label{def:Kbeta}
K_\beta = \Big \{ v: M \to \R \text{ measurable such that }  \beta^{\frac{1}{2}} v = 0 \text{ a.e. in } M \Big \} 
\end{equation}

\begin{equation} \label{eq:def:Ag}
 P_g^s = \Delta_g^s + A_g^s
 \end{equation}

\begin{lemme} \label{lem:mainlemma} Let $s \in \mathbb{N}^*$, $2s <n$, $(M,g)$ be a closed manifold of dimension $n \ge 3$ and $P_g^s$ be the GJMS operator of order $2s$ on $(M,g)$. We assume that \eqref{eq:unique:continuation} holds true. Let $(v_i)_{i \ge0}$ be a sequence in $H^s(M)$, $(\lambda_i)_{i\ge0}$ be a sequence of real numbers, $(p_i)_{i \ge0}$ be a sequence of positive numbers satisfying $p_i \ge \frac{n}{2s}$ for all $i \ge 0$ and, for any $i \ge0$, let $\beta_i \in \mathcal{A}_{p_i}$ be a nonnegative nonzero function. Assume that we have, for all $i \ge 0$, 
\begin{equation}  \label{eq:vecteursvi}
 P_g^s v_i = \lambda_i \beta_i v_i  \quad \text{ in } M, 
 \end{equation}
 $ \int_M \beta_i v_i^2 = 1$, $\int_{M} \beta_i^{p_i} dv_g = 1$, $\lambda_i\neq 0$ and $\limsup_{i\to +\infty} \vert \lambda_i \vert < +\infty$. For all $i \ge 0$ we decompose $v_i$ as 
 $$ v_i = k_i + w_i $$
where $k_i \in \ker(P_g^s)$ and $w_i \in \ker(P_g^s)^{\perp}$. Then 
\begin{enumerate}
\item $ ( w_i )_{i \ge0} $ is bounded in $H^s(M)$
\item There is $c>0$ such that, for all $i\in \mathbb{N}$, $\left\vert \lambda_i \right\vert \geq c$ and $\Vert w_i \Vert_{H^s} \ge c$
\item If $\Vert v_i \Vert_{H^s} \to +\infty$ as $i\to +\infty$ then, $p_i \to \frac{n}{2s}$ from above, $\lambda_i \geq c$ for $i$ large enough, $(w_i)_{i \ge0}$ weakly converges to $0$ in $H^s(M)$, $(\beta_i)_{i \ge0}$ weakly converges to $0$ in $L^q(M)$ for all $q \leq \frac{n}{2s}$ and we have up to the extraction of a subsequence that
$$ v_i = \gamma_i \big(K + o(1) \big) \quad \text{ as } i \to + \infty,$$
where $K$ is a non-zero element of $\ker(P_g^s)$, where $\gamma_i = \Vert v_i \Vert_{H^s}$ and where $o(1)$ denotes a sequence that strongly converges to $0$ in $H^s(M)$.
\end{enumerate}
\end{lemme}

\begin{equation} \label{eq:def:lambdak}
 \lambda_k(\beta)= \inf_{V \in \mathcal{G}_k^\beta(H^s(M))} \max_{v\in V\setminus \{0\}}  \mathcal{R}(\beta,v),
 \end{equation}

\begin{prop} \label{prop:uppersemicontinuity}
Let $p \ge \frac{n}{2s}$ and let  $(\beta_i)_{i\ge0}$ be a sequence of functions in $\mathcal{A}_p$ that converges to some $\beta \in \mathcal{A}_p$ in $L^p(M)$. Then
$$ \limsup_{i\to +\infty} \lambda_k(\beta_i) \leq \lambda_k(\beta) $$
\end{prop}

\begin{equation} \label{eq:enplus:kplus}
%\begin{aligned} 
\lambda_k(\beta)  \ge 0 \quad   \text{ if } k \ge k_+ \text{ and } \beta >0 \text{ a.e.},  \\
\end{equation}

\begin{equation} \label{eq:enplus:kmoins}
\lambda_k(\beta)  \le 0 \quad  \text{ if } k \le k_+-1. 
\end{equation}

\begin{prop} \label{prop:deficontinuityeigen} 
Let $s \in \mathbb{N}^*, 2s <n$ and let $P_g^s$ be the GJMS operator of order $2s$ in $(M,g)$. We assume that $P_g^s$ satisfies \eqref{eq:unique:continuation}. Let $p \ge \frac{n}{2s}$ and $\beta \in \mathcal{A}_p$. We have the following properties
\begin{enumerate}
\item Let $k \in \N^*$, and $b_i \in \mathcal{A}_p$ such that $b_i \to 0 $ as $i\to +\infty$ in $L^p(M)$. Then $\lambda_k(\beta+b_i)\to \lambda_k(\beta)$ in $\mathbb{R} \sqcup \{-\infty\}$ as $i\to +\infty$.
\item There is a sequence of eigenfunctions $\{\vp_k \}_{k\geq m(\beta)}$ in $H^s(M)$ such that for all $k, \ell \ge m(\beta)$ we have $Q(\beta,\vp_k,\vp_\ell) = \delta_{k,\ell}$ and
$$ P_g^s \vp_k = \lambda_k(\beta)\beta \vp_k \quad \text{ in } M. $$
\item We have $\lambda_k(\beta) \to +\infty$ as $k\to +\infty$. For any fixed $k \geq m(\beta)$, the space $E_k(\beta)$ of $k$-th eigenfunctions is equal to
$$ E_k(\beta) = \text{Span} \Big \{  \vp_i ; \lambda_i(\beta) = \lambda_k(\beta) \Big \} $$
and is a finite-dimensional subspace of $H^s(M)$. Furthermore, if there exists $\vp \in H^s(M) \backslash \{0\}$ and $\lambda \in \R$ such that $P_g^s \vp = \lambda \beta \vp$ then $\lambda = \lambda_k(\beta)$ for some $k \ge m(\beta)$ and $\vp \in E_k(\beta)$. 
\item Assume that $m(\beta)\geq 2$. Then for all $m\in \N^*$ and $\epsilon >0$, there are functions $\vp_1^{\epsilon,m},\cdots, \varphi_{m(\beta)-1}^{\epsilon,m}$ that satisfy
 the following: for any $ 1 \leq i,j \leq m(\beta)-1$ and any  $m(\beta) \leq k \leq m$,
$$
 \mathcal{R}(\beta,\vp_i^{\epsilon,m})  \leq -\frac{1}{\epsilon} \text{ and } \left \{
 \begin{aligned}
B(\beta,\vp_i^{\epsilon,m},\vp_j^{\epsilon,m}) & = \delta_{i,j}, \\
\Gamma(\vp_i^{\epsilon,m},\vp_j^{\epsilon,m}) & = \mathcal{R}(\beta,\vp_i^{\epsilon,m}) \delta_{i,j} \\
  B(\beta,\vp_k, \vp_i^{\epsilon,m}) & = 0 \\
 \end{aligned}\right. . $$
% \item Let $k \geq k_+$, and $\beta_i \in \mathcal{A}_p$ such that $\beta_i \to \beta $ as $i\to +\infty$ in $L^p(M)$. Then $\lambda_k(\beta_i)\to \lambda_k(\beta)$ as $i\to +\infty$.
\end{enumerate}
\end{prop}

\begin{prop} \label{prop:wellposedmu_1}
Let $s \in \mathbb{N}^*, 2s <n$ and let $P_g^s$ be the GJMS operator of order $2s$ in $(M,g)$. Let $\beta \in \mathcal{A}_p$ for some $p\geq \frac{n}{2s}$. We assume either that $\beta >0$ a.e or that $P_g^s$ satisfies \eqref{eq:unique:continuation}. Then, the following assumptions are equivalent:
\begin{itemize}
\item[(i)]$ \forall v \in H^s(M)\setminus\{0\}, Q(\beta,v) = 0 \Rightarrow G(v) > 0 $
\item[(ii)] There is $\Lambda \geq 0$ such that if we let, for all $v \in H^s(M)$, %\setminus\{0\}, \quad
$$ N_\beta (v)^2= \Lambda Q(\beta,v) + G(v) ,$$
$N_\beta$ defines a norm on $H^s(M)$ which is equivalent to $\Vert \cdot \Vert_{H^s}$ .
\item[(iii)] $\lambda_1(\beta)>-\infty$.
\end{itemize}
If any of these three assumptions is satisfied %, then $\Vert \cdot \Vert_{H^s}$ and $N_\beta = \sqrt{\Lambda Q(\beta,\cdot) + G(\cdot)}$ are equivalent and
we have
\begin{equation} \label{eq:directsumKbetaHbeta} K_{s,\beta} \oplus H_\beta = H^s(M), \end{equation}
where this splitting is orthogonal for the scalar product associated to $N_{\beta}$, and we have the following min-max characterization of $\lambda_k(\beta)$
\begin{equation} \label{eq:newminmu1} \lambda_k(\beta) = \inf_{V \in \mathcal{G}_k(H_\beta)} \max_{v \in V \setminus \{0\}} \mathcal{R}(\beta,v). \end{equation}
In addition, for any $k \ge 1$, there exists $\vp \in H_\beta$ that satisfies
\begin{equation}\label{eq:eulerlagrangefirsteigen} P_g^s \vp = \lambda_k(\beta) \beta \vp \quad \text{ in } M, \end{equation}
and the vector space $E_k(\beta)$ of solutions of \eqref{eq:eulerlagrangefirsteigen} is finite-dimensional. Finally if, for any $k \ge 1$, $(\vp_{k,\ell}(\beta))_{1 \le \ell \le \dim E_k(\beta)}$ is an orthonormal family of $E_k(\beta)$ with respect to the scalar product associated to $N_{\beta}$ then $\big( \vp_{k,\ell}(\beta)\big)_{k \ge 1, 1 \le \ell \le \dim E_k(\beta)}$ is a Hilbert basis of $H_\beta$ for the same scalar product.
\end{prop}

\begin{proof}
We prove Proposition \ref{prop:deficontinuityeigen} in several steps. 

\medskip

\textbf{Step 1:} let $k \ge 1$. We first claim that $\lambda_k(\beta + \epsilon) \to \lambda_k(\beta)$ as $\epsilon \to 0$. As a consequence, we prove that (2) and (4) hold. 

\medskip 


\textbf{Proof of Step 1:} 
Since $\beta+\ep >0$ everywhere, we have by Proposition \ref{prop:wellposedmu_1} that $\lambda_1(\beta + \ep ) > -\infty$ and that there exits generalised eigenfunctions $(\vp_i^\epsilon)_{i\geq 1}$ associated to $\lambda_i(\beta + \epsilon)$, normalised as 
$$B(\beta+\epsilon,\vp_i^\epsilon,\vp_j^\epsilon) = \delta_{i,j}$$ 
for $i,j \geq 1$.  For simplicity we will let, in the rest of this proof, $\lambda_i^{\epsilon} = \lambda_i(\beta+\epsilon)$ for any $i \ge 1$. Assume first that $k \le m(\beta)-1$. 
Then $\lambda_k(\beta) = -\infty$, and $\lambda_k(\beta + \ep) \to -\infty$ as a consequence of Proposition \ref{prop:uppersemicontinuity}. 
We now assume that $k\ge m(\beta)$. We will prove that there exists $\varphi_k \in H^s(M)$ with $Q(\beta, \varphi_k) = 1$ and $- \infty < \nu_k \le \lambda_k(\beta)$ such that the following holds: 
\begin{equation} \label{eq:eq:contieigen0}
P_g^s \vp_k = \nu_k \beta \vp_k \quad \text{ in } M. 
\end{equation}
Assume first that $k \ge k_+$. Then by \eqref{eq:enplus:kplus} we have $\lambda_k^\ep \geq 0$ for all $\ep >0$, and hence by Proposition \ref{prop:uppersemicontinuity} we obtain $\lambda_k(\beta)\geq 0$. Since $\beta + \epsilon \to \beta$ in $L^p(M)$ as $\epsilon \to 0$ for any $p \ge \frac{n}{2s}$, with $\beta\neq 0$, lemma \ref{lem:mainlemma} applies and shows that the family $(\vp_k^\epsilon)_{\epsilon >0}$ is bounded in $H^s(M)$. Up to the extraction of a subsequence, $\vp_k^\epsilon$ weakly converges  in $H^s(M)$ and strongly converges in $L^2_\beta(M)$ and in $H^{s-1}(M)$ to some function $\vp_k$ as $\ep \to 0$. Up to a subsequence, the previous analysis also shows that $\lambda_k^\epsilon \to \nu_k$ as $\epsilon \to 0$ for some real number $\nu_k$ satisfying $0 < \nu_k \leq \lambda_k(\beta)$. For any $\epsilon >0$, $\vp_k^{\epsilon}$ satisfies 
$$P_g^s \vp_k^{\epsilon} = \lambda_k^{\epsilon}(\beta + \ep) \vp_k^{\epsilon} \quad \text{ in } M, $$
and passing the latter to the weak limit as $\epsilon \to 0$ shows that \eqref{eq:eq:contieigen0} is satisfied when $k \ge k_+$. The relation $Q(\beta+\ve, \vp_k^{\ep} )=1$ also passes to the limit and shows that $Q(\beta, \vp_k) = 1$, hence $\vp_k \neq 0$. If $k, \ell \ge k_+$ the relation $B(\beta+\epsilon,\vp_k^\epsilon,\vp_\ell^\epsilon) = \delta_{k,\ell}$ similarly passes to the limit and shows that $ B(\beta,\vp_k,\vp_\ell) = \delta_{k,\ell}$.  In particular, for $\ep$ small enough,
\begin{equation} \label{eq:dimeigenfunctions}\dim_\beta\left(\text{Span} \left( \vp_i^\ep \right)_{k_+ \leq  i \leq k} \right) = \dim\left(\text{Span} \left( \vp_i^\ep \right)_{k_+ \leq  i \leq k} \right) = k-k_++1. \end{equation}
We now assume that $m(\beta) \leq k\leq k_+-1$. Here the difficulty is to prove that
\begin{equation} \label{eq:provelimsup} \limsup_{\ep\to 0} \vert \lambda_k^\ep \vert < + \infty. \end{equation}
We would like to test $V_\ep = \text{Span}\left(\vp_1^\ep ,\cdots \vp_k^\ep\right)$ in the variational characterization of $\lambda_k(\beta)$. However one could have $\dim_\beta \left( V_\ep\right) < k$. To avoid this problem we set a small perturbation for $t \in \R$:
$$ V_\ep^t =  \text{Span}\left(\vp_i^\ep + t \vp_{k_+ -1 +i}^\ep ; i \in \{1,\cdots,k\} \right). $$
We have that 
$$P(t) := \det\left( \left(\beta^{\frac{1}{2}} \left(\vp_i^\ep + t \vp_{k_+ -1 +i}^\ep \right) \right)_{1\leq i \leq k} \vert \left(\beta^{\frac{1}{2}}\vp_{k_+ -1 + i}^\ep\right)_{1\leq i \leq k} \right)$$
is a well defined polynomial of degree $k$ since $\left(\beta^{\frac{1}{2}}\vp_{k_+ -1 + i}\right)_{1\leq i \leq k}$ is a free family by \eqref{eq:dimeigenfunctions} and we have $\lim_{t\to +\infty} \frac{P(t)}{t^k} = 1$. Since $P$ has finitely many roots we can let $t _\ep \leq 1$ be such that $P(t_\ep)\neq 0$ and we can test $V_\ep^{t_\ep}$ in the definition  \eqref{eq:def:lambdak}  of $\lambda_k(\beta)$.  Let $(\alpha_i^{\ep})_{1 \le i \le k}$ be such that 
$$ v_{\ep}:=  \sum_{i=1}^k \alpha_i^\ep \left(\vp_i^\ep + t_\ep \vp_{k_+ -1 +i}^\ep\right)$$
attains $\max_{v \in V_\ep^{t_\ep} \backslash \{0\}} \frac{G(v)}{Q(\beta,v)} $. Then 
$$ \lambda_k(\beta) \leq \max_{v \in V_\ep^{t_\ep}} \frac{G(v)}{Q(\beta,v)} = \frac{\sum_{i=1}^k \left(\alpha_i^\ep\right)^2\left( \lambda_i^\ep + t_\ep^2 \lambda_{k_+ +i-1}^\ep \right) }{ 1 + t_\ep^2 - Q(\ep,v)} \leq  \frac{\lambda_k^\ep + t_\ep^2 \lambda_{k_+ +k-1}^\ep  }{ 1 + t_\ep^2 - Q(\ep,v)} , $$
%where $v = \sum_{i=1}^k \alpha_i^\ep \left(\vp_i^\ep + t_\ep \vp_{k_+ -1 +i}^\ep\right)$ attains the max with the renormalization $\sum_{i=1}^k \left(\alpha_i^\ep\right)^2 = 1$ and 
where we used that for any $i,j\geq 1$ and $\ep >0$, $\Lambda(\vp_i^\ep,\vp_i^\ep) = \lambda_i^\ep \delta_{i,j}$. Then
$$\lambda_k^\ep \geq \lambda_k(\beta)\left( 1-Q(\ep,v^{\ep}) \right) - t_{\ep}^2 \lambda_{k_++k-1}^\ep \geq \min\left(\lambda_k(\beta),0\right) - \lambda_{k_++k-1}^\ep, $$
where to obtain the last inequality we used \eqref{eq:enplus:kmoins}. The right-hand side of the previous expression is uniformly bounded from below as $\ep \to 0$ since $\limsup_{\ep\to 0} \lambda_{k_++k-1}^\ep \leq \lambda_{k_++k-1}(\beta)$. This proves  \eqref{eq:provelimsup} for $m(\beta) \leq k\leq k_+-1$. With  \eqref{eq:provelimsup} we can proceed as in the case $k\geq k_+$: we apply lemma \ref{lem:mainlemma} to the sequence $(\vp_k^\ep)_{\ep >0}$ and obtain, up to the extraction of a subsequence, that $\vp_k^\ep \rightharpoonup \vp_k$ in $H^s(M)$  and $\lambda_k^\ep \to \nu_k$ for some $-\infty < \nu_k \le \lambda_k(\beta)$ as $\ep \to 0$. Passing the eigenvalue equation satisfied by $\vp_k^{\ep}$ to the weak limit as $\ep \to 0$ similarly proves \eqref{eq:eq:contieigen0} for $k \leq k_+-1$. %Of course we also have up to a subsequence that $\vp_k^\ep \to \vp_k$ in $H^s$ and $\lambda_k^\ep = 0 = \nu_k$ if $k_- +1 \leq k\leq k_+-1$ and $\Delta \vp_k =0$

\smallskip

We now prove (4). For $m \geq m(\beta)$ we let
$$ H_m = \Big \{ v \in H^s(M) ; \forall m(\beta)\leq i \leq m,  Q(\beta,\vp_i,v) = 0  \Big \} $$
where, for $i \ge m(\beta)$, $\vp_i$ is given by \eqref{eq:eq:contieigen0}. We claim that
\begin{equation} \label{eq:mapproximation-infty} \inf_{V \in \mathcal{G}^\beta_{m(\beta)-1}(H_m)} \max_{v\in V\setminus \{0\}} \mathcal{R}(\beta,v) = -\infty .\end{equation}
Indeed, by definition of $m(\beta)$ and by \eqref{eq:def:lambdak} we can let, for any $\epsilon >0$, $V \in \mathcal{G}_{m(\beta)-1}^\beta(H^s(M))$ be such that 
\begin{equation} \label{eq:contieigen1}
 \max_{v\in V\setminus\{0\}} \mathcal{R}(\beta,v) \leq -\frac{1}{\epsilon} 
 \end{equation}
Let
$$ V_m = \left\{ v - \pi_m(v) ; v\in V \right\} $$
where
$$\pi_m(v) = \sum_{i=m(\beta)}^m B(\beta,v,\vp_i) \frac{\vp_i}{Q(\beta,\vp_i)} $$
is the orthogonal projection of $V$ on $ \text{Span} \big \{ \vp_i, m(\beta)\leq i \leq m \big \} $ with respect to $Q(\beta,\cdot)$. We observe that, up to assuming that $\epsilon$ is small enough, we can assume that  $\dim_\beta(V_m) = \dim_\beta(V) = m(\beta)-1$. Indeed, if this were not the case, a linear combination of the family $(\vp_i)_{m(\beta)\leq i \leq m}$ would belong to $V$, which would contradict \eqref{eq:contieigen1} for $\ep$ small enough. Straightforward computations using the definition of $\vp_i$ show that, for any $v \in V$,
$$ B \big( \beta, v- \pi_m(v), \pi_m(v) \big) = \Gamma \big( \beta, v - \pi_m(v), \pi_m(v) \big) = 0$$
holds, so that 
$$ \mathcal{R}(\beta, v - \pi_m(v) ) = \frac{G(v) - G(\pi_m(v)) }{Q(\beta,v) - Q(\beta,\pi_m(v))} \leq \frac{-\frac{1}{\epsilon} Q(\beta,v) - \lambda_{m(\beta)} Q(\beta,\pi_m(v))}{Q(\beta,v)-Q(\beta,\pi_m(v))} \to -\infty$$
as $\epsilon \to 0$. This proves \eqref{eq:mapproximation-infty}. Let now $V_{\epsilon,m} \in \mathcal{G}^\beta_{m(\beta)-1}(H_m)$ be such that 
$$ \max_{v\in V_{\epsilon,m}\setminus \{0\}} \mathcal{R}(\beta,v) \leq - \frac{1}{\epsilon} $$
By an orthodiagonalization of $\Gamma$ with respect to $Q(\beta,\cdot)$ on $V_{\epsilon,m}$, we obtain a family $\vp_1^{\epsilon,m},\cdots,\vp_{m(\beta)-1}^{\epsilon,m}$ which satisfies property (4).

\smallskip 

We are now in position to conclude the proof of (1) when $b_i = \epsilon_i$ and  $\epsilon_i \to 0$ is a sequence of positive constants, and conclude the proof of (2).  Let $k \ge m(\beta)$ and let
$$V_{\epsilon}  = \text{Span} \Big \{ \vp_1^{\epsilon,k},\cdots,\vp_{m(\beta)-1}^{\epsilon,k},\vp_{m(\beta)},\cdots, \vp_{k} \Big \}. $$
The variational characterization \eqref{eq:def:lambdak} of $\lambda_k(\beta)$ together with \eqref{eq:eq:contieigen0} show that 
\begin{equation} \label{eq:contieigen2} \begin{aligned}
\lambda_k(\beta) & \leq \max_{\alpha \in \mathbb{S}^{k-1}} \frac{\sum_{i=1}^{m(\beta)-1}\alpha_i^2 \mathcal{R}(\beta,  \vp_i^{\epsilon,m})+  \sum_{i=m(\beta)}^k \alpha_i^2 \mathcal{R}(\beta,  \vp_i)}{\sum_{i=1}^{m(\beta)-1}\alpha_i^2 Q(\beta,  \vp_i^{\epsilon,m})+  \sum_{i=m(\beta)}^k \alpha_i^2 Q(\beta,  \vp_i)} \\
& \le  \max_{\alpha \in \mathbb{S}^{k-1}}  \sum_{i=m(\beta)}^m \alpha_i^2 \mathcal{R}(\beta,  \vp_i) \leq \nu_k.
\end{aligned} 
\end{equation}
Therefore, $\nu_k = \lambda_k(\beta)$ and (1) holds for $b_i = \ep_i$. Coming back to \eqref{eq:eq:contieigen0}, this also finishes the proof of point (2). 


\medskip

\textbf{Step 2:} We prove property (3).

\medskip


\textbf{Proof of Step 2:} First, recall that by Proposition \ref{prop:wellposedmu_1} $E_k(\beta)$ is orthogonal to $K_{\beta}$ defined in \eqref{def:Kbeta} for the bilinear form $\Gamma(\beta, \cdot, \cdot)$. As a consequence, $B(\beta, \cdot, \cdot)$ defines a scalar product on $E_k(\beta)$. We first prove that the unit ball of $E_k(\beta)$ with respect to the norm $\sqrt{Q(\beta,\cdot)}$ is compact. Indeed, let $(v_i)_{i \ge0}$ be a sequence of generalised eigenfunctions associated to $\lambda_k(\beta)$. By Lemma \ref{lem:mainlemma}, since $\beta \neq 0$ is fixed, $(v_i)_{i \ge0}$ is bounded in $H^s(M)$. Up to the extraction of a subsequence, $v_i$ converges to some $v$ weakly in $H^s(M)$ and strongly in $L^2_{\beta}(M)$. Passing to the weak limit, it is easily seen that $v$ still satisfies $P_g^s v = \lambda_k(\beta) \beta v$ in $M$. For all $i \ge 0$ we thus have $P_g^s(v_i - v) = \lambda_k(\beta) \beta (v_i-v)$, so that integrating against $v_i-v$ and using \eqref{eq:def:Ag} shows that $ \int_M \big|\Delta_g^{\frac{s}{2}} (v_i-v) \big|^2 dv_g = o(1)$  as $i \to + \infty$. Thus the unit ball of  $E_k(\beta)$ for $\sqrt{Q(\beta,\cdot)}$ is compact and $E_k(\beta)$ is finite dimensional. 

Assume now by contradiction that the sequence $(\lambda_k(\beta))_{k\ge m(\beta)}$ is bounded. A similar argument to the one we just used, again relying on Lemma \ref{lem:mainlemma}, then shows that the unit ball of $E = \overline{\bigoplus_{k\geq m(\beta)} E_k(\beta)}$ with respect to the norm $\sqrt{Q(\beta,\cdot)}$ is compact. This shows that $E$ is finite-dimensional, which is impossible since the family $(\vp_k)_{k\ge m(\beta)}$ is free by point (2) of Proposition  \ref{prop:deficontinuityeigen}. Hence $\lim_{k \to +\infty} \lambda_k(\beta) = + \infty$. 

Assume now that there are $\vp\in H^s(M) \backslash \{0\}$ and $\lambda\in \R$ such that $P_g^s \vp = \lambda \beta \vp$ in $M$ and assume that $\lambda \not \in \{\lambda_k(\beta)\}_{k \ge m(\beta)}$. We let $k \in \N^*$ be such that $\lambda_{k-1}(\beta)<  \lambda < \lambda_k(\beta) $. We let 
$$ \tilde{H}_k = \Big \{ v \in H^s(M) ; Q(\beta, \vp, v) = 0 \text{ and}, \forall m(\beta)\leq i \leq k-1,  Q(\beta,\vp_i,v) = 0  \Big \} .$$
Arguing as in the proof of point (4) above it is easily seen that, for any $\epsilon >0$, there exists a family $\tilde{\vp}_1^{\epsilon}, \cdots, \tilde{\vp}_{m(\beta)-1}^{\epsilon}$ of functions in $H^s(M)$ satisfying, for any $ 1 \leq i,j \leq m(\beta)-1$ and any  $m(\beta) \leq \ell \leq k-1$,
$$
 \mathcal{R}(\beta,\tilde{\vp}_i^{\epsilon})  \leq -\frac{1}{\epsilon} \text{ and } \left \{
 \begin{aligned}
B(\beta,\tilde{\vp}_i^{\epsilon},\tilde{\vp}_j^{\epsilon}) & = \delta_{i,j}, \\
\Gamma(\tilde{\vp}_i^{\epsilon},\tilde{\vp}_j^{\epsilon}) & = \mathcal{R}(\beta,\tilde{\vp}_i^{\epsilon}) \delta_{i,j} \\
  B(\beta,\vp_\ell, \tilde{\vp}_i^{\epsilon}) & = 0 \\
  B(\beta, \vp,  \tilde{\vp}_i^{\epsilon})& = 0
 \end{aligned}\right. . $$
We let
$$V_{\epsilon}  = \text{Span} \Big \{ \tilde{\vp}_1^{\epsilon},\cdots,\tilde{\vp}_{m(\beta)-1}^{\epsilon},\vp_{m(\beta)},\cdots, \vp \Big \}. $$
Then $\dim_{\beta} V_\ep = k$ and mimicking the computations in \eqref{eq:contieigen2}, the variational characterisation \eqref{eq:def:lambdak} of $\lambda_k(\beta)$ then shows that $\lambda_k(\beta) \le \lambda$, which is a contradiction. Hence $\lambda = \lambda_m(\beta)$ for some $m \ge m(\beta)$, and this also proves that for all $k \ge m(\beta)$ we have 
$$ E_k(\beta) = \text{Span} \Big \{  \vp_i ; \lambda_i(\beta) = \lambda_k(\beta) \Big \}, $$
and concludes the proof of (3). 

\medskip

\textbf{Step 3:} we conclude the proof of (1). 

\medskip 

\textbf{Proof of Step 3:} we let $(b_i)_{i \ge 0}$ be a sequence in $\mathcal{A}_p$ such that $b_i \to 0$ in $L^p(M)$ as $ i \to + \infty$. We claim that 
\begin{equation} \label{semicontinuite:m}
m(\beta+ b_i) \le m(\beta)
\end{equation} 
for all $i$ large enough. Assume first that $\lambda_{m(\beta)}(\beta) \ge 0$. Then, by \eqref{eq:def:lambdak}, any $V \in \mathcal{G}_k^{\beta}(H^s(M))$ contains a nonzero element $v$ such that $G(v) \ge 0$. Then, since $b_i \ge 0$, 
$$\max_{v \in V \backslash \{0\}} \frac{G(v)}{Q(\beta + b_i, v)} \ge0$$ 
and hence $\lambda_{m(\beta)}(\beta+b_i) \ge 0$. Assume now that $\lambda_{m(\beta)} <0$. Proposition~\ref{prop:uppersemicontinuity} then shows that $\lambda_{m(\beta)}(\beta + b_i) < 0$ for $i$ large enough. In the definition \eqref{eq:def:lambdak} of $\lambda_k(\beta + b_i)$ the infimum may therefore only be taken over subsets $V \in \mathcal{G}_k^{\beta}(H^s(M))$ such that $\max_{v \in V \backslash \{0\}} \mathcal{R}(\beta, v)  <0$. For a such a $V$ we have, since $b_i \ge 0$, 
$$\max_{v \in V \backslash \{0\}} \frac{G(v)}{Q(\beta + b_i, v)} \ge \max_{v \in V \backslash \{0\}} \frac{G(v)}{Q(\beta,v)},$$
and taking the infimum over $V$ yields $\lambda_{m(\beta)}(\beta + b_i) \ge \lambda_{m(\beta)}(\beta)$. This proves \eqref{semicontinuite:m}. 
%Assume first that $m(\beta) = 1$.  XXXXX . Assume now that $m(\beta) \ge 2$. We let $\vp_1^{\epsilon},\cdots, \varphi_{m(\beta)-1}^{\epsilon}$ be the functions given by  point (4) of Proposition~\ref{prop:deficontinuityeigen} applied with $m = m(\beta)$ and $\epsilon >0$. They satisfy 
%% and  there are functions $\vp_1^{\epsilon,m},\cdots, \varphi_{m(\beta)-1}^{\epsilon,m}$ that satisfy
% the following: for any $ 1 \leq i,j \leq m(\beta)-1$,
%$$
% \mathcal{R}(\beta,\vp_i^{\epsilon,m})  \leq -\frac{1}{\epsilon} \text{ and } \left \{
% \begin{aligned}
%B(\beta,\vp_i^{\epsilon,m},\vp_j^{\epsilon,m}) & = \delta_{i,j}, \\
%\Gamma(\vp_i^{\epsilon,m},\vp_j^{\epsilon,m}) & = \mathcal{R}(\beta,\vp_i^{\epsilon,m}) \delta_{i,j} \\
%  B(\beta,\vp_{m(\beta)}, \vp_i^{\epsilon,m}) & = 0 \\
% \end{aligned}\right. . $$
% Let $V \in \mathcal{G}_k^\beta(H^s(M))$.
 % \inf_{V \in \mathcal{G}_k^\beta(H^s(M))} \max_{v\in V\setminus \{0\}}  \mathcal{R}(\beta,v),
With \eqref{semicontinuite:m}, and by points (2) and (3) of Proposition~\ref{prop:deficontinuityeigen}, for any $i \ge 0$ there thus exists a sequence $(\varphi_k^i)_{k \ge m(\beta)}$ of generalised eigenfunctions associated to $\beta + b_i$ and satisfying $P_g^s \vp_k^i = \lambda_k(\beta + b_i) (\beta + b_i) \vp_k^i$ in $M$ for every $k \ge m(\beta)$ and $i \ge 0$. The arguments in the proof of \eqref{semicontinuite:m} shows that for every $k \ge m(\beta)$ the sequence $(\lambda_k(\beta + b_i))_{i \ge 0}$ is uniformly bounded and, up to a subsequence, converges to some $\nu_k \in \R$ with $\nu_k \le \lambda_k(\beta)$. We can now mimic the arguments in the proof of Step $1$ to show that $( \vp_k^i)_{i \ge0}$ converges, up to a subsequence, to $\vp_k \in H^s(M)$ satisfying $P_g^s \vp_k = \nu_k \beta \vp_k$ in $M$. That $\nu_k = \lambda_k(\beta)$ now follows from point (3) of Proposition~\ref{prop:deficontinuityeigen}.
\end{proof}

\end{document}
